\section{二维随机变量及其分布规律}

本节介绍二维随机变量及其的分布律。

本节要点：
\begin{itemize}
    \item 掌握二维离散型随机变量的及其分布律的概念；
    \item 掌握二维连续型随机变量的及其密度函数的概念。
\end{itemize}

%============================================================
\subsection{二维离散型随机变量及其分布律的概念}

\begin{definition}[二维离散型随机变量及其分布律]
若二维随机变量$\left( X,Y \right) $的取值为有限个或可列无穷多个，则称$\left( X,Y \right) $为{\bf 二维离散型随机变量}，简称{\bf 二维离散变量}。
设所有可能的取值设为$\left( x_i,y_j \right) ,i=1,2,\cdots ,j=1,2,\cdots $，则对应的概率记作：
\[
P\left\{ X=x_i,Y=y_j \right\} =p_{ij}
\]
称为{\bf 二维离散型随机变量$\left( X,Y \right) $的分布律}，一般可用二维表格罗列。
其中：
\begin{itemize}
    \item $x_i$：纯数学记号，对应事件$A_i$；
    \item $y_i$：纯数学记号，对应事件$B_i$；
    \item $p_{ij}=P\left\{ X=x_i,Y=y_j \right\} =P\left( A_iB_j \right) $：$A_i,B_j$同时发生的概率。
\end{itemize}
\end{definition}

若二维离散变量$\left( X,Y \right) $有分布律$p_{ij}$，则其分布函数为：
\[
F\left( x,y \right) =P\left\{ X\leqslant x,Y\leqslant y \right\} =\sum_{x_i\leqslant x,y_j\leqslant y}{p_{ij}}
\]
由分布函数的充要条件可得$p_{ij}$成为分布律的充要条件：
\[
p_{ij}\text{为分布律}\Leftrightarrow \begin{cases}
	p_{ij}\geqslant 0\\
	\sum_j{\sum_i{p_{ij}}}=1\\
\end{cases}
\]
其中：
\begin{itemize}
    \item $\sum_j{\sum_i{p_{ij}}}$：表示将所有的$p_{ij}$累加起来。
\end{itemize}

\begin{tcolorbox}
再次注意，分布律$p_{i,j}$表示两个事件都发生。
\end{tcolorbox}

%============================================================
\subsection{二维连续型随机变量及其密度函数的概念}

\begin{definition}[二维连续型随机变量及其密度函数]
设二维随机变量$\left( X,Y \right) $及其分布函数$F\left( x,y \right) $，若存在二元可积函数$f\left( x,y \right) \geqslant 0$，对于$\forall x,y$，有：
\[
F\left( x,y \right) =\iint\limits_D{f\left( u,v \right) \cdot d\sigma},\qquad \begin{array}{l}
	d\sigma =dudv\\
	D=\left\{ \left( u,v \right) \middle| u\leqslant x,v\leqslant y \right\}\\
\end{array}
\]
则称$\left( X,Y \right) $为{\bf 二维连续型随机变量}，简称{\bf 二维连续变量}，并称$f\left( x,y \right) $为{\bf 二维连续型随机变量$\left( X,Y \right) $的密度函数}。
\end{definition}

由分布函数的充要条件易得$f\left( x,y \right) $为密度函数的充要条件：
\[
f\left( x,y \right) \text{为密度函数}\Leftrightarrow \begin{cases}
	f\left( x,y \right) \geqslant 0\\
	\iint\limits_{\mathbb{R} \times \mathbb{R}}{f\left( x,y \right) \cdot d\sigma}=1,\quad d\sigma =dxdy\\
\end{cases}
\]
对于任意平面$D$的概率有：
\[
P\left\{ \left( X,Y \right) \in D \right\} =\iint\limits_D{f\left( x,y \right) \cdot d\sigma},\qquad d\sigma =dxdy
\]

\begin{tcolorbox}
同样，密度函数$f\left( x,y \right) $表示两个事件都发生的概率密度。
其次，一个二维变量是否为连续型，看它有没有密度函数，至于变量本身是整个连续还是分段连续并不要紧。
\end{tcolorbox}


\begin{theorem}
设二维连续随机变量$\left( X,Y \right) $有分布函数$F\left( x,y \right) $和密度函数$f\left( x,y \right) $，若$f\left( x,y \right) $在点$\left( x_0,y_0 \right) $处连续，则：
\[
f\left( x_0,y_0 \right) =\frac{\partial ^2F\left( x_0,y_0 \right)}{\partial x\partial y}=\frac{\partial ^2F\left( x_0,y_0 \right)}{\partial y\partial x}
\]
\end{theorem}

可用微积分的二阶混偏相等定理证明，略。

%============================================================
\subsection{例}

\begin{example}
若5个产品中有2个次品，每次抽取一个，取两次，记$X$表示第一次取到次品的个数，$Y$表示第二次取到次品的个数，试求：
\begin{enumerate}
    \item 不放回抽取的情况下，$\left( X,Y \right) $的分布律，
    \item 放回抽取的情况下，$\left( X,Y \right) $的分布律。
\end{enumerate}
\end{example}

解：

由题干可得，$X,Y$取值只可能是0或1，0表示抽到合格品，1表示抽到次品。

1. 不放回抽取的情况下：
\begin{align*}
&P\left\{ X=0,Y=0 \right\} =\frac{3}{5}\cdot \frac{2}{4}=\frac{3}{10} \\
&P\left\{ X=1,Y=0 \right\} =\frac{2}{5}\cdot \frac{3}{4}=\frac{3}{10} \\
&P\left\{ X=0,Y=1 \right\} =\frac{3}{5}\cdot \frac{2}{4}=\frac{3}{10} \\
&P\left\{ X=1,Y=1 \right\} =\frac{2}{5}\cdot \frac{1}{4}=\frac{1}{10}
\end{align*}

2. 放回抽取的情况下：
\begin{align*}
&P\left\{ X=0,Y=0 \right\} =\frac{3}{5}\cdot \frac{3}{5}=\frac{9}{25} \\
&P\left\{ X=1,Y=0 \right\} =\frac{2}{5}\cdot \frac{3}{5}=\frac{6}{25} \\
&P\left\{ X=0,Y=1 \right\} =\frac{3}{5}\cdot \frac{2}{5}=\frac{6}{25} \\
&P\left\{ X=1,Y=1 \right\} =\frac{2}{5}\cdot \frac{2}{5}=\frac{4}{25}
\end{align*}

~

\begin{example}
若二维连续随机变量$\left( X,Y \right) $有密度函数
\[
f\left( x,y \right) =\begin{cases}
	ke^{-x},&		y\in \left( 0,x \right)\\
	0,&		y\notin \left( 0,x \right)\\
\end{cases}
\]
计算分布函数和$P\left\{ X+Y<2 \right\} $。
\end{example}

解：

首先得计算具体的密度函数，由密度函数的充要条件得：
\begin{align*}
&\because \iint\limits_D{f\left( x,y \right) \cdot d\sigma}=1,\quad D=\left\{ \left( x,y \right) \middle| x>0,0<y<x \right\} \\
&\because \iint\limits_D{ke^{-x}\cdot d\sigma}=\int_0^{+\infty}{dx}\int_0^x{ke^{-x}\cdot dy}=\int_0^{+\infty}{kxe^{-x}\cdot dx}=k \\
&\therefore k=1 \\
&\therefore f\left( x,y \right) =\begin{cases}
	e^{-x},&		y\in \left( 0,x \right)\\
	0&		y\notin \left( 0,x \right)\\
\end{cases}
\end{align*}
使用了部分积分法$\int{xe^xdx}=\int{xd\left( e^x \right)}=xe^x-\int{e^xdx}$。

\begin{figure}[ht]
	\centering
	\includegraphics[height=2.5cm]{3.2.png}
	\end{figure}

$P\left\{ X+Y<2 \right\} $的区域为$\left\{ \left( x,y \right) \middle| x\in \left[ 0,2 \right] ,y\in \left[ 0,\min \left\{ x,-x+2 \right\} \right] \right\} $
。由于积分区域是凸区域，所以积分结果和次序无关：
\[
P\left\{ \left( X,Y \right) \in D \right\} =\iint\limits_D{e^{-x}\cdot d\sigma}=\int_0^1{dy}\int_y^{2-y}{e^{-x}\cdot dx}=1-2e^{-1}+e^{-2}
\]




